% ============================================================
% 论文摘要 PDF 统一样式模板（配色 + 风格 + 完整 preamble）
% 用法：在 paper.tex 中 \documentclass[11pt,a4paper]{ctexart} 之后
%       % ============================================================
% 论文摘要 PDF 统一样式模板（配色 + 风格 + 完整 preamble）
% 用法：在 paper.tex 中 \documentclass[11pt,a4paper]{ctexart} 之后
%       % ============================================================
% 论文摘要 PDF 统一样式模板（配色 + 风格 + 完整 preamble）
% 用法：在 paper.tex 中 \documentclass[11pt,a4paper]{ctexart} 之后
%       % ============================================================
% 论文摘要 PDF 统一样式模板（配色 + 风格 + 完整 preamble）
% 用法：在 paper.tex 中 \documentclass[11pt,a4paper]{ctexart} 之后
%       \input{assets/paper_style.tex}
% 仅含 preamble，不含 \documentclass / \begin{document}
% 配色与组件统一，勿在正文自创新颜色或新环境
% ============================================================

% --- 数学 ---
\usepackage{amsmath,amssymb,amsthm,mathtools,bm}
% --- 图表 ---
\usepackage{graphicx,booktabs,array,multirow,makecell}
\usepackage[table]{xcolor}
% --- 排版 ---
\usepackage[margin=2.3cm]{geometry}
\usepackage{enumitem}        % 紧凑列表
\usepackage{tcolorbox}       % 提示框/分析框
\tcbuselibrary{skins,breakable}
\usepackage{titlesec}
\usepackage[colorlinks=true,linkcolor=teal,citecolor=teal,
            bookmarksnumbered=true]{hyperref}  % 书签目录便于 PDF 导航

% --- 配色（统一，勿随意新增） ---
\definecolor{brandblue}{HTML}{2E54C8}
\definecolor{okgreen}{HTML}{2E7D32}
\definecolor{warnorange}{HTML}{EF6C00}
\definecolor{badred}{HTML}{C62828}
\definecolor{deeppurple}{HTML}{6A1B9A}
\definecolor{headgray}{HTML}{283593}

% --- 章节标题配色（headgray） ---
\titleformat{\section}{\Large\bfseries\color{headgray}}{\thesection}{0.6em}{}
\titleformat{\subsection}{\large\bfseries\color{headgray}}{\thesubsection}{0.6em}{}
\titleformat{\subsubsection}{\normalsize\bfseries\color{headgray}}{\thesubsubsection}{0.6em}{}

% --- 提示框：统一可复用环境，颜色按内容性质选 ---
% 用法：\begin{callout}[okgreen] ... \end{callout}（缺省 brandblue）
\newtcolorbox{callout}[1][brandblue]{
  colback=#1!6, colframe=#1, boxrule=0.8pt, arc=2pt,
  left=8pt,right=8pt,top=6pt,bottom=6pt, breakable}

% --- 数值例子框：用于"喂进去什么、出来什么"的实例化解释 ---
% 用法：\begin{example} ... \end{example}
\newtcolorbox{example}{
  enhanced, colback=gray!4, colframe=headgray!70, boxrule=0.6pt, arc=2pt,
  left=8pt,right=8pt,top=6pt,bottom=6pt, breakable,
  borderline west={3pt}{0pt}{headgray},
  fonttitle=\bfseries, title={数值例子}}

% --- 定理环境 ---
\newtheorem{theorem}{定理}
\newtheorem{corollary}{推论}

% 仅含 preamble，不含 \documentclass / \begin{document}
% 配色与组件统一，勿在正文自创新颜色或新环境
% ============================================================

% --- 数学 ---
\usepackage{amsmath,amssymb,amsthm,mathtools,bm}
% --- 图表 ---
\usepackage{graphicx,booktabs,array,multirow,makecell}
\usepackage[table]{xcolor}
% --- 排版 ---
\usepackage[margin=2.3cm]{geometry}
\usepackage{enumitem}        % 紧凑列表
\usepackage{tcolorbox}       % 提示框/分析框
\tcbuselibrary{skins,breakable}
\usepackage{titlesec}
\usepackage[colorlinks=true,linkcolor=teal,citecolor=teal,
            bookmarksnumbered=true]{hyperref}  % 书签目录便于 PDF 导航

% --- 配色（统一，勿随意新增） ---
\definecolor{brandblue}{HTML}{2E54C8}
\definecolor{okgreen}{HTML}{2E7D32}
\definecolor{warnorange}{HTML}{EF6C00}
\definecolor{badred}{HTML}{C62828}
\definecolor{deeppurple}{HTML}{6A1B9A}
\definecolor{headgray}{HTML}{283593}

% --- 章节标题配色（headgray） ---
\titleformat{\section}{\Large\bfseries\color{headgray}}{\thesection}{0.6em}{}
\titleformat{\subsection}{\large\bfseries\color{headgray}}{\thesubsection}{0.6em}{}
\titleformat{\subsubsection}{\normalsize\bfseries\color{headgray}}{\thesubsubsection}{0.6em}{}

% --- 提示框：统一可复用环境，颜色按内容性质选 ---
% 用法：\begin{callout}[okgreen] ... \end{callout}（缺省 brandblue）
\newtcolorbox{callout}[1][brandblue]{
  colback=#1!6, colframe=#1, boxrule=0.8pt, arc=2pt,
  left=8pt,right=8pt,top=6pt,bottom=6pt, breakable}

% --- 数值例子框：用于"喂进去什么、出来什么"的实例化解释 ---
% 用法：\begin{example} ... \end{example}
\newtcolorbox{example}{
  enhanced, colback=gray!4, colframe=headgray!70, boxrule=0.6pt, arc=2pt,
  left=8pt,right=8pt,top=6pt,bottom=6pt, breakable,
  borderline west={3pt}{0pt}{headgray},
  fonttitle=\bfseries, title={数值例子}}

% --- 定理环境 ---
\newtheorem{theorem}{定理}
\newtheorem{corollary}{推论}

% 仅含 preamble，不含 \documentclass / \begin{document}
% 配色与组件统一，勿在正文自创新颜色或新环境
% ============================================================

% --- 数学 ---
\usepackage{amsmath,amssymb,amsthm,mathtools,bm}
% --- 图表 ---
\usepackage{graphicx,booktabs,array,multirow,makecell}
\usepackage[table]{xcolor}
% --- 排版 ---
\usepackage[margin=2.3cm]{geometry}
\usepackage{enumitem}        % 紧凑列表
\usepackage{tcolorbox}       % 提示框/分析框
\tcbuselibrary{skins,breakable}
\usepackage{titlesec}
\usepackage[colorlinks=true,linkcolor=teal,citecolor=teal,
            bookmarksnumbered=true]{hyperref}  % 书签目录便于 PDF 导航

% --- 配色（统一，勿随意新增） ---
\definecolor{brandblue}{HTML}{2E54C8}
\definecolor{okgreen}{HTML}{2E7D32}
\definecolor{warnorange}{HTML}{EF6C00}
\definecolor{badred}{HTML}{C62828}
\definecolor{deeppurple}{HTML}{6A1B9A}
\definecolor{headgray}{HTML}{283593}

% --- 章节标题配色（headgray） ---
\titleformat{\section}{\Large\bfseries\color{headgray}}{\thesection}{0.6em}{}
\titleformat{\subsection}{\large\bfseries\color{headgray}}{\thesubsection}{0.6em}{}
\titleformat{\subsubsection}{\normalsize\bfseries\color{headgray}}{\thesubsubsection}{0.6em}{}

% --- 提示框：统一可复用环境，颜色按内容性质选 ---
% 用法：\begin{callout}[okgreen] ... \end{callout}（缺省 brandblue）
\newtcolorbox{callout}[1][brandblue]{
  colback=#1!6, colframe=#1, boxrule=0.8pt, arc=2pt,
  left=8pt,right=8pt,top=6pt,bottom=6pt, breakable}

% --- 数值例子框：用于"喂进去什么、出来什么"的实例化解释 ---
% 用法：\begin{example} ... \end{example}
\newtcolorbox{example}{
  enhanced, colback=gray!4, colframe=headgray!70, boxrule=0.6pt, arc=2pt,
  left=8pt,right=8pt,top=6pt,bottom=6pt, breakable,
  borderline west={3pt}{0pt}{headgray},
  fonttitle=\bfseries, title={数值例子}}

% --- 定理环境 ---
\newtheorem{theorem}{定理}
\newtheorem{corollary}{推论}

% 仅含 preamble，不含 \documentclass / \begin{document}
% 配色与组件统一，勿在正文自创新颜色或新环境
% ============================================================

% --- 数学 ---
\usepackage{amsmath,amssymb,amsthm,mathtools,bm}
% --- 图表 ---
\usepackage{graphicx,booktabs,array,multirow,makecell}
\usepackage[table]{xcolor}
% --- 排版 ---
\usepackage[margin=2.3cm]{geometry}
\usepackage{enumitem}        % 紧凑列表
\usepackage{tcolorbox}       % 提示框/分析框
\tcbuselibrary{skins,breakable}
\usepackage{titlesec}
\usepackage[colorlinks=true,linkcolor=teal,citecolor=teal,
            bookmarksnumbered=true]{hyperref}  % 书签目录便于 PDF 导航

% --- 配色（统一，勿随意新增） ---
\definecolor{brandblue}{HTML}{2E54C8}
\definecolor{okgreen}{HTML}{2E7D32}
\definecolor{warnorange}{HTML}{EF6C00}
\definecolor{badred}{HTML}{C62828}
\definecolor{deeppurple}{HTML}{6A1B9A}
\definecolor{headgray}{HTML}{283593}

% --- 章节标题配色（headgray） ---
\titleformat{\section}{\Large\bfseries\color{headgray}}{\thesection}{0.6em}{}
\titleformat{\subsection}{\large\bfseries\color{headgray}}{\thesubsection}{0.6em}{}
\titleformat{\subsubsection}{\normalsize\bfseries\color{headgray}}{\thesubsubsection}{0.6em}{}

% --- 提示框：统一可复用环境，颜色按内容性质选 ---
% 用法：\begin{callout}[okgreen] ... \end{callout}（缺省 brandblue）
\newtcolorbox{callout}[1][brandblue]{
  colback=#1!6, colframe=#1, boxrule=0.8pt, arc=2pt,
  left=8pt,right=8pt,top=6pt,bottom=6pt, breakable}

% --- 数值例子框：用于"喂进去什么、出来什么"的实例化解释 ---
% 用法：\begin{example} ... \end{example}
\newtcolorbox{example}{
  enhanced, colback=gray!4, colframe=headgray!70, boxrule=0.6pt, arc=2pt,
  left=8pt,right=8pt,top=6pt,bottom=6pt, breakable,
  borderline west={3pt}{0pt}{headgray},
  fonttitle=\bfseries, title={数值例子}}

% --- 定理环境 ---
\newtheorem{theorem}{定理}
\newtheorem{corollary}{推论}
